\documentclass{article}
\usepackage{amsmath, amssymb}

\title{Theorem 1: Exponential Stability of the Flux Diagnostic}
\author{Dionisio Alberto Lopez III}
\date{January 30, 2026}

\begin{document}

\maketitle

\section{Theorem Statement}

\textbf{Theorem 1 (Nekhoroshev Confinement Indicator)}. 
For a near-integrable Hamiltonian system $H(I, \theta) = H_0(I) + \varepsilon H_1(I, \theta)$ with confinement functional $Q(I)$, if the system satisfies the Nekhoroshev steepness conditions and the perturbation parameter $\varepsilon$ is below the threshold $\varepsilon_0$, then the time-averaged flux diagnostic satisfies:

\begin{equation}
|\langle Z_Q \rangle_T| < C \cdot \exp\left(-\frac{\kappa}{\varepsilon^\alpha}\right)
\end{equation}

for all times $T < T_{Nekh} \sim \exp(\varepsilon_0/\varepsilon)^{1/2n}$, where $n$ is the number of degrees of freedom.

\section{Proof Outline}

\subsection{1. Bounds on Action Drift}
From Nekhoroshev's theorem (1977), we have the bound on action variation:
\begin{equation}
|I(t) - I(0)| < \varepsilon^a \quad \text{for} \quad t < T_{Nekh}
\end{equation}

\subsection{2. Relation to Confinement Functional}
Let $Q$ be a smooth function of the actions $I$. The time derivative $\dot{Q}$ is given by:
\begin{equation}
\dot{Q} = \nabla_I Q \cdot \dot{I} = \nabla_I Q \cdot \{I, H\}
\end{equation}
Since $\dot{I} \sim \mathcal{O}(\varepsilon)$, $\dot{Q}$ is also small.

\subsection{3. Time Averaging}
The diagnostic is $Z_Q = Q \dot{Q} / E$. The time average is:
\begin{equation}
\langle Z_Q \rangle_T = \frac{1}{T} \int_0^T \frac{Q(t) \dot{Q}(t)}{E} dt = \frac{1}{TE} \left[ \frac{1}{2} Q(t)^2 \right]_0^T
\end{equation}
Since $Q(t)$ is bounded by the Nekhoroshev estimate, the term in brackets is bounded. As $T \to \infty$ (or $T \sim T_{Nekh}$), the average decays.
However, the strictly exponential bound comes from the fact that the \textit{secular} drift component of $\dot{Q}$ is exponentially small in the stable regime.

\end{document}
